% Sources: README "Adaptive Run Length", "Spectrum Thinning"; PLAN.md section 7

\section{Fitting method}
\label{sec:fitting}

\subsection{The sources}
\label{sec:fitting:sources}
% Source: literature table from an ADS search (papers after 2014) made in another session, passed by the author
%         2026-11-13; bibcodes in refs.bib (all VERIFY)

The sample consists of the 14 ULXs of Table~\ref{tab:sources}, named in the text and the tables by their code names.
In four of them pulsations have been detected, so the accretor is a NS.
For the others no pulsations are known; we call them BH candidates where the cited works interpret their spectral and timing properties as super-Eddington accretion onto a BH, and NS candidates where the cited works favour a NS.
Neither kind is a confirmed identification.
% AUTHOR: the code names come from the SED files (SED_keV_ergps_<name>.txt); the SED script
%         (../data/scripts/ISSI_ULX_make_SED_topsources.sh) labels NGC1313B, NGC1313H, NGC5408, NGC247S, NGC247SS as
%         "NGC1313 BD", "NGC1313 HUL", "NGC5408 SUL", "NGC247 SUL", "NGC247 SSUL". If these are spectral regimes
%         (broadened disc, hard/soft/supersoft ultraluminous; Claude's reading), the suffixes name states, not sources:
%         NGC1313B and NGC1313H may both be NGC 1313 X-1 (then NGC1313H is not the pulsar X-2), and NGC247S/SS the
%         same source in two states. "ULS" in M51ULS, M81ULS, M101ULS may mean ultraluminous supersoft sources, which
%         would not be M51 ULX-7 (a hard PULX) or M81 X-6. Identifications marked "?" in the table: please confirm.

\begin{table*}
\centering
\caption{The sources. PULX: pulsating ULX. ``?'': identification of the code name not confirmed.}
\label{tab:sources}
\footnotesize
\begin{tabular}{@{}lp{0.17\textwidth}p{0.1\textwidth}p{0.33\textwidth}p{0.2\textwidth}@{}}
\toprule
Code name & Identification & Type & Evidence & References \\
\midrule
HolmII & Holmberg~II X-1 & BH candidate & no pulsations; spectra consistent with super-Eddington accretion onto a BH; NS not excluded & \citealt{Barra2024, ReyeroSerantes2024} \\
HolmIX & Holmberg~IX X-1 & BH candidate & no pulsations; a new spectral state; treated as a BH ULX & \citealt{Walton2025} \\
M101ULS & M101 ULX-1 & BH candidate & BH mass ${\sim}20\,M_\odot$ from spectral scaling; no pulsations & \citealt{Titarchuk2016} \\
M51ULS & M51 ULX-7? & NS (PULX) & $P_{\rm spin} = 2.8$~s, $P_{\rm orb} = 2$~d & \citealt{RodriguezCastillo2020} \\
M81ULS & M81 X-6? & NS candidate & spectral and timing analysis favours a NS; no pulsations & \citealt{Amato2023} \\
NGC1313B & NGC~1313 X-1? & BH candidate & no pulsations; classified as a BH ULX & \citealt{Majumder2023, Majumder2025} \\
NGC1313H & NGC~1313 X-2? & NS (PULX) & $P_{\rm spin} \approx 1.5$~s & \citealt{Sathyaprakash2019} \\
NGC247S & NGC~247 ULX-1 & NS candidate & supersoft ULX; NS favoured by several groups & \citealt{Pinto2021, Zhou2023} \\
NGC247SS & NGC~247 ULX-1? & NS candidate & the same source as NGC247S in its supersoft state; no paper on a second ULX in NGC~247 & \citealt{DAi2021, Alston2021} \\
NGC300 & NGC~300 ULX-1 (SN~2010da) & NS (PULX) & $P_{\rm spin} = 31.6$~s, rapid spin-up & \citealt{Carpano2018} \\
NGC5408 & NGC~5408 X-1 & BH candidate & QPOs, spectral behaviour of a BH ULX; no pulsations & \citealt{Majumder2023} \\
NGC55 & NGC~55 ULX-1 & BH candidate & classified as a BH ULX in long-term spectral studies; no confirmed pulsations & \citealt{Majumder2025, Jithesh2022} \\
NGC6946 & NGC~6946 X-1 & BH candidate & two-temperature disc spectra; classified as a BH ULX & \citealt{Ghosh2023, Majumder2023} \\
NGC7793 & NGC~7793 P13 & NS (PULX) & $P_{\rm spin} = 0.42$~s & \citealt{Furst2016, Israel2017} \\
\bottomrule
\end{tabular}
% source: as above (author 2026-11-13). Left out of the evidence column: "He donor star" for NGC 247 ULX-1 (not
%         checked; the other session's own note says its NGC 55 pulsation claim was not found in ADS).
\end{table*}

\subsection{Where the SEDs come from}
\label{sec:fitting:data}

The SEDs that we fit are not observed spectra but a model fitted to them.
The X-ray spectrum of each source was fitted in SPEX over 0.3--10~keV only with an absorbed modified blackbody plus Comptonization, $\code{hot} \times (\code{mbb} + \code{comt})$, with the seed temperature of \code{comt} tied to the temperature of \code{mbb}.
The absorbing column was then set to zero, and the unabsorbed model was evaluated from $10^{-4}$ to 100~keV.
Below 0.3~keV and above 10~keV the SED is therefore the extrapolation of this model, and its slope there tells nothing about the disc.
Its low-energy slope, ${\rm d}\log\nu L_\nu/{\rm d}\log E \to 1.75$, is the asymptote of the SPEX modified blackbody \citep{KaastraBarr1989}, whose photon spectrum is $N(E) \propto E^{1/4}/[{\rm e}^{x}({\rm e}^{x} - 1)]^{1/2}$ with $x = E/kT$ (see also Sect.~\ref{sec:disc:thermalization}).
% source: ULXspectra-mc docs/thermalization_in_wind.md sect. 1 (commit 69b9de7); checked (Claude): x << 1 gives N ~ E^(-1/4), nuLnu ~ E^(7/4)
% source: ../data/scripts/ISSI_ULX_make_SED_topsources.sh (SPEX: ign 0:0.3, ign 10:100; par 1 3 t0 c 1 2 t;
%         par 1 1 fc v 0; plot rx 1e-4 1e2 -> SED_keV_ergps_*.txt); author 2026-11-04
% AUTHOR: who made the SEDs and from which observations (instrument, ObsID); cite SPEX? (not in refs.bib)
% OUTLINE (plan, author 2026-11-04):
%   - fits restricted to 0.3-10 keV (the data) vs the full range (mostly extrapolation); the "full"
%     fits weigh the extrapolated mbb and comt tails in chi^2 -> say what each fit measures;
%   - ask the coauthors whether real UV/optical counterpart fluxes (HST, Swift/UVOT, XMM/OM) and
%     hard X-ray data (NuSTAR) exist for these sources, and why mbb rather than a disc component;
%   - NGC1313B / NGC1313H: the script labels them "NGC1313 BD" and "NGC1313 HUL" (spectral states);
%     identify the source(s) for the names.

% OUTLINE: Data: SEDs of the selected ULXs (data/selected); LMIN = 1e36 erg/s cut;
%          full energy range vs <= 1 keV (see sec:fitting:data above).
% OUTLINE: Likelihood in log nuLnu with sigma_log = 0.08; 20 points per decade of
%          energy (equal weight per decade) and why (SEDs smooth; Fisher test).
% OUTLINE: MCMC: emcee with differential-evolution moves (ter Braak 2006; ter Braak &
%          Vrugt 2008), 200 walkers, adaptive burn-in and run length (N >= 20 tau,
%          ESS ~ 4000).
